\scanpage{34}
\begin{center}$-\ 29\ -$\end{center}
\hfill n$^\circ$ 371

\begin{center}\S\ 2. \uuline{Produits tensoriels d'algèbres.}\end{center}

1. \uline{Produit tensoriel d'un nombre fini d'algèbres.}

\uline{Dans ce n$^\circ$, nous nous écartons de nos conventions générales, en entendant par ``algèbre'' une algèbre associative mais n'ayant pas nécessairement d'élément unité.}

Soient $A$ un anneau commutatif, $(E_i)_{i\in I}$ une famille finie de $A$-\uline{algèbres} (associatives, mais n'ayant pas nécessairement d'élément unité). Soit $E=\bigotimes_{i\in I}E_i$ le $A$-\uline{module} produit tensoriel des $A$-modules $E_i$ (chap. II, $3^{\mathrm e}$ éd., \S\ 3, n$^\circ$ 9). La structure d'algèbre de $E_i$ est définie par une application $A$-\uline{linéaire} $m_i:E_i\otimes E_i\to E_i$ (\S\ 1, n$^\circ$ 1, \uline{Remarque} 1) ; on déduit par suite de ces applications une application $A$-linéaire
\[
m=\bigotimes_i m_i:\bigotimes_{i\in I}(E_i\otimes E_i)\longrightarrow\bigotimes_{i\in I}E_i=E
\]
et en vertu de l'associativité du produit tensoriel, cette application s'identifie à une application $A$-linéaire de $E\otimes E$ dans $E$ ou encore à une application $A$-bilinéaire $E\times E\to E$ ; cette dernière définit une structure de $A$-\uline{algèbre} (associative, mais n'ayant pas nécessairement d'élément unité) sur $E$ ; en effet elle est telle, par définition, que l'on ait
\[
(1)\qquad \left(\bigotimes_i x_i\right)\left(\bigotimes_i y_i\right)=\bigotimes_i(x_i y_i)
\]
et par suite on a bien
\[
(2)\qquad \left(\left(\bigotimes_i x_i\right)\left(\bigotimes_i y_i\right)\right)\left(\bigotimes_i z_i\right)=\bigotimes_i(x_i y_i z_i)=\left(\bigotimes_i x_i\right)\left(\left(\bigotimes_i y_i\right)\left(\bigotimes_i z_i\right)\right)
\]
en vertu de l'associativité des $E_i$ ; la relation $(xy)z=x(yz)$ dans $E$ en résulte par linéarité.

DEFINITION 1.--- \uline{Etant donnée une famille finie} $(E_i)_{i\in I}$ \uline{d'algèbres (n'ayant pas nécessairement d'élément unité) sur un anneau commutatif} $A$, \uline{on appelle produit}
\mspage{33}
